\documentclass[12pt]{article}
\usepackage{fullpage}
\usepackage[T1]{fontenc}
\usepackage[utf8]{inputenc}
\usepackage{amsmath,amssymb,amsthm}
\usepackage{hyperref}

\newcommand{\GL}{\mathrm{GL}}
\newcommand{\oo}{\mathfrak{o}}
\newcommand{\pp}{\mathfrak{p}}
\newcommand{\qq}{\mathfrak{q}}
\newcommand{\Wmod}{\mathcal{W}}
\newcommand{\diag}{\operatorname{diag}}
\DeclareMathOperator{\Cc}{C_c^\infty}

\newtheorem{theorem}{Theorem}
\newtheorem{lemma}{Lemma}
\newtheorem{proposition}{Proposition}

\title{Problem 2: A uniform test vector for the local Rankin--Selberg integral}
\author{}
\date{}

\begin{document}
\maketitle

\begin{theorem}
The answer is \emph{yes}. Let $\Pi$ be a generic irreducible admissible
representation of $\GL_{n+1}(F)$, realized in $\Wmod(\Pi,\psi^{-1})$. Then there
exists $W\in\Wmod(\Pi,\psi^{-1})$ such that for every generic irreducible
admissible representation $\pi$ of $\GL_n(F)$ there is a vector
$V\in\Wmod(\pi,\psi)$ for which the local Rankin--Selberg integral
\[
  \Psi(s,W,V):=\int_{N_n\backslash \GL_n(F)} W(\diag(g,1)\,u_Q)\,V(g)\,|\det g|^{s-\frac12}\,dg
\]
is finite and nonzero for all $s\in\mathbb C$. (Here, as in the statement,
$\qq$ is the conductor ideal of $\pi$, $Q$ generates $\qq^{-1}$, and
$u_Q=I_{n+1}+Q\,E_{n,n+1}$.)
\end{theorem}

Throughout, $F$ is a non-archimedean local field with ring of integers $\oo$,
maximal ideal $\pp$, and residue field of size $q_F$; $\psi$ has conductor $\oo$,
i.e.\ $\psi|_{\oo}=1$ and $\psi|_{\pp^{-1}}\neq1$. We write $K=\GL_n(\oo)$ for the
maximal compact subgroup, $N_r\subset\GL_r$ for the upper triangular unipotent
subgroup, $P_r\subset\GL_r$ for the mirabolic subgroup (matrices whose bottom row
is $(0,\dots,0,1)$), and $e_r=(0,\dots,0,1)^{\mathsf T}$. We use throughout that a
function on a homogeneous space is ``$(N_r,\chi)$-equivariant'' if
$f(ng)=\chi(n)f(g)$ for $n\in N_r$, and $\Cc(N_r\backslash G,\chi)$ denotes the
smooth such functions with support compact modulo $N_r$.

The proof has three steps.
\begin{itemize}
\item[(A)] We use the mirabolic (Bernstein--Zelevinsky/Kirillov) structure of
$\Pi$ to choose, \emph{once and for all}, a single vector $W$ depending only on
$\Pi$. We compute $\Psi(s,W,V)$ explicitly and find that it collapses to an
$s$-independent linear functional $L(V)$ of $V$.
\item[(B)] We prove a prescription lemma for Whittaker functions of $\pi$: their
restrictions to any finite union of $N_n$--$K'$ double cosets are completely
arbitrary (subject only to equivariance).
\item[(C)] We deduce that $L$ is not identically zero on $\Wmod(\pi,\psi)$ by a
positivity (Gauss-sum) argument, producing the required $V$.
\end{itemize}

\section{The fixed vector $W$ and the collapse of the integral}

\subsection{The mirabolic radical and the effect of $u_Q$}
\label{sec:mirabolic}

Embed $\GL_n\hookrightarrow\GL_{n+1}$ by $g\mapsto\diag(g,1)$. Both $\diag(g,1)$
and $u_Q=I_{n+1}+Q\,E_{n,n+1}$ lie in $P_{n+1}$: indeed the bottom row of $u_Q$ is
$e_{n+1}^{\mathsf T}=(0,\dots,0,1)$. Hence $\diag(g,1)u_Q\in P_{n+1}$ for all
$g\in\GL_n$.

Identify $P_{n+1}=\GL_n\ltimes F^n$ via
\[
  (h,v)\longleftrightarrow\begin{pmatrix} h & v\\ 0 & 1\end{pmatrix},
  \qquad h\in\GL_n,\ v\in F^n\ (\text{column}),
\]
with multiplication $(h,v)(h',v')=(hh',hv'+v)$. In these coordinates
$\diag(g,1)=(g,0)$, and since $u_Q E_{n,n+1}$ places $Q$ times the $n$-th column of
$\diag(g,1)$, namely $Q\,g e_n$, into the last column,
\begin{equation}\label{eq:uQ}
  \diag(g,1)\,u_Q=(g,\;Q\,g e_n),\qquad g e_n=\text{$n$-th column of }g.
\end{equation}

The subgroup $N_{n+1}\subset P_{n+1}$ consists of $(u,y)$ with $u\in N_n$ and
$y\in F^n$ arbitrary; its generic character $\psi$ evaluates on $(u,y)$ as
$\psi$ of the standard sum of superdiagonal entries of $u$, times $\psi(y_n)$
(the $(n,n+1)$-entry of $\diag(u,1)\cdot(I,y)$ is $y_n$). Thus
$\psi^{-1}((I,y))=\psi(-y_n)$.

\begin{lemma}\label{lem:Cc-iso}
The restriction map $f\mapsto F$, $F(g):=f((g,0))$, is a linear isomorphism
\[
  \Cc(N_{n+1}\backslash P_{n+1},\psi^{-1})\;\xrightarrow{\ \sim\ }\;
  \Cc(N_n\backslash \GL_n,\psi^{-1}),
\]
and for every $f$ in the source and every $g\in\GL_n$,
\begin{equation}\label{eq:twist}
  f\bigl(\diag(g,1)\,u_Q\bigr)=\psi(-Q\,g_{nn})\,F(g),
\end{equation}
where $g_{nn}$ is the $(n,n)$-entry of $g$.
\end{lemma}

\begin{proof}
For $(h,v)\in P_{n+1}$ we have $(h,v)=(I,v)(h,0)$, and $(I,v)\in N_{n+1}$; left
$(N_{n+1},\psi^{-1})$-equivariance of $f$ gives
$f((h,v))=\psi^{-1}((I,v))\,f((h,0))=\psi(-v_n)\,F(h)$.
Thus $f$ is determined by $F$. The function $F$ is left $(N_n,\psi^{-1})$-equivariant
because for $u\in N_n$, $(u,0)\in N_{n+1}$ and
$F(ug)=f((ug,0))=f((u,0)(g,0))=\psi^{-1}((u,0))F(g)=\psi^{-1}(u)F(g)$; here
$\psi^{-1}((u,0))$ is the standard character of $u\in N_n$. Smoothness and the
compact-support condition correspond exactly under the identification
$N_{n+1}\backslash P_{n+1}=N_n\backslash\GL_n$, $(h,0)\mapsto h$. Conversely any
$F\in\Cc(N_n\backslash\GL_n,\psi^{-1})$ defines $f((h,v)):=\psi(-v_n)F(h)$, which
is well defined and lies in the source (the consistency
$f((u,0)(h,v))=\psi^{-1}(u)f((h,v))$ holds by the equivariance of $F$). This proves
the isomorphism.

For \eqref{eq:twist}, use \eqref{eq:uQ}: $\diag(g,1)u_Q=(g,Q\,ge_n)$, and the
last coordinate of $Q\,g e_n$ is $Q\,(g e_n)_n=Q\,g_{nn}$. Hence
$f((g,Q\,ge_n))=\psi(-Q\,g_{nn})F(g)$.
\end{proof}

\subsection{The chosen vector $W$}
\label{sec:chosenW}

By the theory of the mirabolic subgroup for generic representations
(Bernstein--Zelevinsky derivatives; equivalently the Kirillov model of $\Pi$
restricted to $P_{n+1}$), the restriction map
$W\mapsto W|_{P_{n+1}}$ from $\Wmod(\Pi,\psi^{-1})$ to functions on $P_{n+1}$ is
injective and its image \emph{contains} $\Cc(N_{n+1}\backslash P_{n+1},\psi^{-1})$.
(This is the statement that the ``top derivative'' part
$\Phi^+(\cdots)\Psi^+$ of the restriction of a generic representation to the
mirabolic is the compactly induced representation
$\Cc(N_{n+1}\backslash P_{n+1},\psi^{-1})$; see
Bernstein--Zelevinsky, \emph{Ann.\ Sci.\ ENS} 10 (1977), \S5, and
Jacquet--Piatetski-Shapiro--Shalika, \emph{Amer.\ J.\ Math.} 105 (1983).)

\paragraph{Definition of $W$.}
Let $F_0\in\Cc(N_n\backslash\GL_n,\psi^{-1})$ be the function supported on $N_nK$
and given by
\begin{equation}\label{eq:F0}
  F_0(uk)=\psi^{-1}(u)\qquad(u\in N_n,\ k\in K),\qquad F_0\equiv0\ \text{off}\ N_nK.
\end{equation}
This is well defined: if $uk=u'k'$ with $u,u'\in N_n$, $k,k'\in K$, then
$u'^{-1}u=k'k^{-1}\in N_n\cap K=N_n(\oo)$, and since the superdiagonal entries of
$N_n(\oo)$ lie in $\oo$ where $\psi=1$, we get $\psi^{-1}(u)=\psi^{-1}(u')$. Also
$F_0$ is right $K$-invariant and supported on the set $N_nK$, which is compact
modulo $N_n$ (as $(N_n\cap K)\backslash K$ is compact); hence
$F_0\in\Cc(N_n\backslash\GL_n,\psi^{-1})$. Note $F_0$ depends only on $F$, not on
$\pi$.

Let $f_0\in\Cc(N_{n+1}\backslash P_{n+1},\psi^{-1})$ be the corresponding function
under Lemma~\ref{lem:Cc-iso}, and choose $W\in\Wmod(\Pi,\psi^{-1})$ with
$W|_{P_{n+1}}=f_0$ (possible by the displayed mirabolic surjectivity). \emph{This
$W$ depends only on $\Pi$.}

\subsection{Computation of the integral}

\begin{proposition}\label{prop:collapse}
For $W$ as above and any $V\in\Wmod(\pi,\psi)$,
\begin{equation}\label{eq:collapse}
  \Psi(s,W,V)=c_0\cdot L(V),\qquad c_0:=\operatorname{vol}\bigl((N_n\cap K)\backslash N_nK\bigr)>0,\qquad
  L(V):=\int_{(N_n\cap K)\backslash K}\psi(-Q\,k_{nn})\,V(k)\,dk,
\end{equation}
which is independent of $s$. Consequently $\Psi(s,W,V)$ is finite for every $s$,
and it is nonzero for all $s\in\mathbb C$ if and only if $L(V)\neq0$.
\end{proposition}

\begin{proof}
By the definition of $W$ and \eqref{eq:twist},
$W(\diag(g,1)u_Q)=f_0(\diag(g,1)u_Q)=\psi(-Q\,g_{nn})F_0(g)$. Since $F_0$ is
supported on $N_nK$, the integrand vanishes unless $g\in N_nK$. Choose a Haar
measure on $N_n\backslash\GL_n$; over the compact set $N_n\backslash N_nK$ we may
integrate over a fundamental domain $(N_n\cap K)\backslash K$ for $N_n$ acting on
$N_nK$. Writing $g=uk$ with $u\in N_n$, $k\in K$ (representative $k\in K$),
we have:
\begin{itemize}
\item $F_0(uk)=\psi^{-1}(u)$ by \eqref{eq:F0};
\item $V(uk)=\psi(u)V(k)$ by $(N_n,\psi)$-equivariance of $V$;
\item $(uk)_{nn}=k_{nn}$, because the last row of $u\in N_n$ is $e_n^{\mathsf T}$,
so the last row of $uk$ equals the last row of $k$;
\item $|\det g|=|\det k|=1$, since $k\in K$.
\end{itemize}
Hence the integrand on $N_n\backslash N_nK$ equals
$\psi(-Q\,k_{nn})\,\psi^{-1}(u)\,\psi(u)\,V(k)\cdot1=\psi(-Q\,k_{nn})V(k)$, which is
independent of $u$ and of $s$. Integrating over the fundamental domain yields a
positive constant (the volume of the $N_n$-fibre, absorbed into the normalization)
times $L(V)$. As $(N_n\cap K)\backslash K$ is compact and the integrand is locally
constant, $L(V)$ is a finite complex number, and $\Psi(s,W,V)$ equals this number
for every $s$. The final assertion is immediate.
\end{proof}

We have reduced everything to:
\begin{equation}\label{eq:goal}
  \boxed{\ \exists\,V\in\Wmod(\pi,\psi)\ \text{with}\ L(V)\neq0,\quad
  L(V)=\int_{(N_n\cap K)\backslash K}\psi(-Q\,k_{nn})\,V(k)\,dk.\ }
\end{equation}

\section{A prescription lemma for Whittaker functions}

We will need the freedom to choose $V$ with prescribed restriction to $K$. The
following is the key local statement; it makes precise the heuristic ``the
Whittaker model is locally as large as possible.''

\begin{lemma}[Single-coset prescription]\label{lem:single}
Let $\pi$ be generic irreducible admissible on $\GL_n(F)$, with Whittaker model
$\Wmod(\pi,\psi)$. Let $g_0\in\GL_n$ and let $K'\subseteq\GL_n$ be a compact open
subgroup. Then there is a vector $V\in\Wmod(\pi,\psi)$ that is right $K'$-invariant,
supported (modulo $N_n$) inside $N_n g_0 K'$, and satisfies $V(g_0)=1$, provided
$K'$ is small enough (depending on $g_0$).
\end{lemma}

\begin{proof}
By the Kirillov model of $\pi$ (Bernstein--Zelevinsky; Jacquet--Piatetski-Shapiro--Shalika),
the restriction map $V\mapsto V|_{P_n}$ is injective and its image contains
$\Cc(N_n\backslash P_n,\psi)$; in particular for any $\phi_0\in\Cc(N_n\backslash
P_n,\psi)$ there is $V_0\in\Wmod(\pi,\psi)$ with $V_0|_{P_n}=\phi_0$.

The group $\GL_n$ acts transitively on nonzero row vectors $\xi\in F^n\setminus\{0\}$
by $g\mapsto e_n^{\mathsf T}g$, with stabilizer of $e_n^{\mathsf T}$ equal to $P_n$;
thus $P_n\backslash\GL_n\cong F^n\setminus\{0\}$. Pick $y_0\in\GL_n$ with the same
bottom row as $g_0$; then $p_0:=g_0 y_0^{-1}\in P_n$.

Choose $\phi_0\in\Cc(N_n\backslash P_n,\psi)$ supported in a small $N_n$-neighbourhood
$N_n\,\Omega_0$ of $p_0$ (with $\Omega_0\subseteq P_n$ a small compact open set,
$p_0\in\Omega_0$) and with $\phi_0(p_0)=1$. Let $V_0\in\Wmod(\pi,\psi)$ realize
$\phi_0$ on $P_n$, and put $V_1:=\pi(y_0)V_0$, i.e.\ $V_1(g)=V_0(gy_0^{-1})$. Then
\[
  V_1(g_0)=V_0(g_0 y_0^{-1})=V_0(p_0)=\phi_0(p_0)=1,
\]
and $V_1\in\Wmod(\pi,\psi)$ (the model is stable under right translation). Because
$V_1$ is smooth, it is right invariant under some compact open subgroup $K''$; and
because $\phi_0$ has small support, $V_1$ is supported (mod $N_n$) in a small
neighbourhood of $N_n g_0$. Shrinking $K'$ so that $K'\subseteq K''$ and so that
$N_n g_0 K'$ contains the support of $V_1$ (mod $N_n$), the vector $V:=V_1$ has all
required properties. The smallness of $K'$ depends only on $g_0$ (through $y_0$ and
the support of $V_1$). 
\end{proof}

\begin{lemma}[Finite-set prescription]\label{lem:finite}
Let $K'=1+\pp^{m}M_n(\oo)$ be a principal congruence subgroup with $m\ge1$. Let
$\Theta:N_n\backslash\GL_n\to\mathbb C$ be left $(N_n,\psi)$-equivariant, right
$K'$-invariant, and supported on finitely many $N_n$--$K'$ double cosets
$N_n g_1 K',\dots,N_n g_M K'$ with $g_1,\dots,g_M\in K$. If $m$ is large enough
(depending on $g_1,\dots,g_M$), then there is $V\in\Wmod(\pi,\psi)$ with
$V|_{N_n g_j K'}=\Theta|_{N_n g_j K'}$ for all $j$.
\end{lemma}

\begin{proof}
By Lemma~\ref{lem:single} applied to each $g_j$ (taking $m$ large enough for all
$j$), there are $V_j\in\Wmod(\pi,\psi)$, right $K'$-invariant, supported (mod $N_n$)
in $N_n g_j K'$, with $V_j(g_j)=1$. Since $\Theta$ is right $K'$-invariant and
$(N_n,\psi)$-equivariant, on the coset $N_n g_j K'$ it is determined by its value
$\Theta(g_j)$ (any two elements of $N_n g_j K'$ representing the same point of
$N_n\backslash\GL_n$ are related on the left by $N_n$ and on the right by $K'$, and
both $\Theta$ and $V_j$ transform by the same rules). Hence
$\Theta|_{N_n g_j K'}=\Theta(g_j)\,V_j|_{N_n g_j K'}$. Because the cosets
$N_n g_j K'$ are pairwise disjoint, the vector $V:=\sum_{j=1}^{M}\Theta(g_j)\,V_j$
satisfies $V|_{N_n g_j K'}=\Theta|_{N_n g_j K'}$ for each $j$, and is supported on
their union, as required. (Consistency of the per-coset assignment is automatic
because each $V_j$ is itself an honest $(N_n,\psi)$-equivariant vector.)
\end{proof}

\section{Non-vanishing of $L$ and conclusion}

Let $c\ge0$ be the conductor exponent of $\pi$, so $\qq=\pp^{c}$ and
$Q\in\qq^{-1}=\pp^{-c}$ with $v(Q)=-c$. Because $\psi$ has conductor $\oo$,
\begin{equation}\label{eq:cond}
  \psi(-Qx)=1\iff x\in\pp^{c}\qquad(x\in F).
\end{equation}

\begin{lemma}\label{lem:phi}
The function $\varphi:K\to\mathbb C$, $\varphi(k):=\psi(-Q\,k_{nn})$, is right
invariant under the principal congruence subgroup $K_c:=1+\pp^{c}M_n(\oo)$ and is
left invariant under $N_n\cap K$ (equivalently, $\varphi$ defines a left
$(N_n,\psi)$-equivariant function on $N_nK$ via $uk\mapsto\psi(u)\varphi(k)$).
\end{lemma}

\begin{proof}
\emph{Right $K_c$-invariance.} For $\kappa\in K_c$ write $\kappa=1+\beta$ with
$\beta\in\pp^{c}M_n(\oo)$. Then
$(k\kappa)_{nn}=\sum_{\ell}k_{n\ell}\kappa_{\ell n}
=k_{nn}\kappa_{nn}+\sum_{\ell\ne n}k_{n\ell}\kappa_{\ell n}$.
Here $\kappa_{nn}\in1+\pp^{c}$ and $\kappa_{\ell n}\in\pp^{c}$ for $\ell\ne n$, while
$k_{n\ell}\in\oo$; hence
$(k\kappa)_{nn}=k_{nn}+\bigl(k_{nn}(\kappa_{nn}-1)+\sum_{\ell\ne n}k_{n\ell}\kappa_{\ell n}\bigr)
\in k_{nn}+\pp^{c}$. By \eqref{eq:cond},
$\psi(-Q(k\kappa)_{nn})=\psi(-Qk_{nn})$.

\emph{Left $(N_n\cap K)$-compatibility.} For $u\in N_n\cap K=N_n(\oo)$ the last row
of $u$ is $e_n^{\mathsf T}$, so $(uk)_{nn}=k_{nn}$ and $\varphi(uk)=\varphi(k)$; on
the other hand $\psi(u)=1$ since the superdiagonal entries of $u$ lie in $\oo$.
Hence $uk\mapsto\psi(u)\varphi(k)=\varphi(uk)$ is a well-defined left
$(N_n,\psi)$-equivariant function on $N_nK$.
\end{proof}

\begin{proposition}\label{prop:nonvanish}
There exists $V\in\Wmod(\pi,\psi)$ with $L(V)\neq0$.
\end{proposition}

\begin{proof}
Take $K'=K_c$ if $c\ge1$, and $K'=1+\pp\,M_n(\oo)$ if $c=0$; in either case $K'$ is
a principal congruence subgroup, and by Lemma~\ref{lem:phi} the function
$\varphi(k)=\psi(-Qk_{nn})$ is right $K'$-invariant. Define
$\Theta:N_nK\to\mathbb C$ by
\[
  \Theta(uk):=\psi(u)\,\overline{\varphi(k)}=\psi(u)\,\psi(Q\,k_{nn})
  \qquad(u\in N_n,\ k\in K),
\]
and $\Theta\equiv0$ off $N_nK$. By Lemma~\ref{lem:phi} (applied to
$\overline\varphi=\psi(Q\,\cdot\,)$, which has the same invariance properties),
$\Theta$ is a well-defined left $(N_n,\psi)$-equivariant, right $K'$-invariant
function. Its support $N_nK$ meets only finitely many $N_n$--$K'$ double cosets,
since $(N_n\cap K)\backslash K/K'$ is finite ($K/K'$ is finite); choose
representatives $g_1,\dots,g_M\in K$. Enlarging the level $m$ of $K'$ if necessary
(replacing $K'$ by a smaller principal congruence subgroup on which $\varphi$ and
$\Theta$ are still invariant — this only refines the partition into cosets), we may
assume $m$ is large enough for Lemma~\ref{lem:finite}. That lemma then provides
$V\in\Wmod(\pi,\psi)$ with $V|_{N_n g_j K'}=\Theta|_{N_n g_j K'}$ for all $j$; thus
$V|_{N_nK}=\Theta$, i.e.\ $V(k)=\psi(Q\,k_{nn})$ for $k\in K$.

For this $V$,
\[
  L(V)=\int_{(N_n\cap K)\backslash K}\psi(-Q\,k_{nn})\,V(k)\,dk
     =\int_{(N_n\cap K)\backslash K}\psi(-Q\,k_{nn})\,\psi(Q\,k_{nn})\,dk
     =\int_{(N_n\cap K)\backslash K}1\,dk
     =\operatorname{vol}\bigl((N_n\cap K)\backslash K\bigr)>0.
\]
Both factors $\psi(-Qk_{nn})$ and $V(k)=\psi(Qk_{nn})$ are constant on
$N_n$-cosets (the first by Lemma~\ref{lem:phi}, the second by
$(N_n,\psi)$-equivariance and $\psi|_{N_n\cap K}=1$), so the integrand descends to
the quotient and equals $1$. Hence $L(V)\neq0$.
\end{proof}

\begin{proof}[Proof of the Theorem]
Fix $\Pi$ and let $W\in\Wmod(\Pi,\psi^{-1})$ be the vector constructed in
\S\ref{sec:chosenW}; it depends only on $\Pi$. Let $\pi$ be an arbitrary generic
irreducible admissible representation of $\GL_n(F)$, with conductor ideal $\qq$ and
$Q$ a generator of $\qq^{-1}$. By Proposition~\ref{prop:nonvanish} there is
$V\in\Wmod(\pi,\psi)$ with $L(V)\neq0$. By Proposition~\ref{prop:collapse},
$\Psi(s,W,V)$ is independent of $s$ and equals a positive constant times $L(V)$,
hence is finite and nonzero for every $s\in\mathbb C$. This is exactly the asserted
property of $W$.
\end{proof}

\section*{Remarks}

\begin{enumerate}
\item The mechanism is the following. The shift $u_Q$ converts, via
\eqref{eq:twist}, the Rankin--Selberg integrand into $\psi(-Q\,g_{nn})$ times a
\emph{fixed} compactly supported function $F_0$ coming from the mirabolic model of
$\Pi$. The single coordinate $g_{nn}$ is the only one touched by the twist; on the
maximal compact $K$, where $|\det|\equiv1$, the integral becomes the $s$-independent
oscillatory integral $L(V)=\int_K\psi(-Q\,k_{nn})V(k)\,dk$. The exponent $-c$ of $Q$
matches the conductor of $\pi$, so $\psi(-Q\,\cdot\,)$ is exactly of level $c$; the
vector $V$ may be chosen of the same depth to cancel the oscillation, after which
$L(V)$ is a (positive) integral of $1$. This is the same phenomenon that makes
local Gauss sums (epsilon factors) nonzero.

\item Because the resulting integral is constant in $s$, it is in particular a
nonzero element of $\mathbb C[q_F^{\pm s}]$ with neither zeros nor poles, which is
the precise meaning of ``finite and nonzero for all $s\in\mathbb C$.'' (A bare
$L$-factor would not work, as it has poles.)

\item No genericity beyond admissibility/genericity of $\pi$ and $\Pi$ is used; the
construction is uniform in $\pi$ once $W$ is fixed.
\end{enumerate}

\end{document}
